% !TeX program = xelatex
\documentclass[UTF8,openany]{ctexbook}

% ---------- 页面与基础功能 ----------
\usepackage[a4paper,centering,body={16cm,25cm}]{geometry}
\usepackage{graphicx}
\usepackage{float}
\usepackage[dvipsnames,svgnames]{xcolor}
\usepackage{hyperref}
\usepackage{bookmark}
\usepackage{enumitem}
\usepackage{siunitx}
\usepackage{import}

% ---------- 数学 ----------
\usepackage{amsmath,amssymb,mathtools,amsthm,bm,mathrsfs}
\usepackage{physics2}
\usephysicsmodule{ab}

% 兼容旧式 Dirac 记号：\braket{\psi|\phi}
\usepackage{braket}
\let\bra\Bra
\let\ket\Ket
\let\braket\Braket

% ---------- 定理框 ----------
\usepackage{mdframed}

% ---------- Quarto / Pandoc 正文兼容 ----------
% Quarto / Pandoc 生成正文所需的最小兼容层。
% 这里只放正文可能调用的命令与环境，不包含 documentclass、标题页或目录。

% ---------- 代码块 ----------
\usepackage{fancyvrb}
\usepackage{framed}

\newcommand{\VerbBar}{|}
\newcommand{\VERB}{\Verb[commandchars=\\\{\}]}
\DefineVerbatimEnvironment{Highlighting}{Verbatim}{commandchars=\\\{\}}

\definecolor{shadecolor}{RGB}{241,243,245}
\newenvironment{Shaded}{\begin{snugshade}}{\end{snugshade}}

\newcommand{\AlertTok}[1]{\textcolor[rgb]{0.68,0.00,0.00}{#1}}
\newcommand{\AnnotationTok}[1]{\textcolor[rgb]{0.37,0.37,0.37}{\textit{#1}}}
\newcommand{\AttributeTok}[1]{\textcolor[rgb]{0.40,0.45,0.13}{#1}}
\newcommand{\BaseNTok}[1]{\textcolor[rgb]{0.68,0.00,0.00}{#1}}
\newcommand{\BuiltInTok}[1]{\textcolor[rgb]{0.00,0.23,0.31}{#1}}
\newcommand{\CharTok}[1]{\textcolor[rgb]{0.13,0.47,0.30}{#1}}
\newcommand{\CommentTok}[1]{\textcolor[rgb]{0.37,0.37,0.37}{#1}}
\newcommand{\CommentVarTok}[1]{\textcolor[rgb]{0.37,0.37,0.37}{\textit{#1}}}
\newcommand{\ConstantTok}[1]{\textcolor[rgb]{0.56,0.35,0.01}{#1}}
\newcommand{\ControlFlowTok}[1]{\textcolor[rgb]{0.00,0.23,0.31}{\textbf{#1}}}
\newcommand{\DataTypeTok}[1]{\textcolor[rgb]{0.68,0.00,0.00}{#1}}
\newcommand{\DecValTok}[1]{\textcolor[rgb]{0.68,0.00,0.00}{#1}}
\newcommand{\DocumentationTok}[1]{\textcolor[rgb]{0.37,0.37,0.37}{\textit{#1}}}
\newcommand{\ErrorTok}[1]{\textcolor[rgb]{0.68,0.00,0.00}{#1}}
\newcommand{\ExtensionTok}[1]{\textcolor[rgb]{0.00,0.23,0.31}{#1}}
\newcommand{\FloatTok}[1]{\textcolor[rgb]{0.68,0.00,0.00}{#1}}
\newcommand{\FunctionTok}[1]{\textcolor[rgb]{0.28,0.35,0.67}{#1}}
\newcommand{\ImportTok}[1]{\textcolor[rgb]{0.00,0.46,0.62}{#1}}
\newcommand{\InformationTok}[1]{\textcolor[rgb]{0.37,0.37,0.37}{#1}}
\newcommand{\KeywordTok}[1]{\textcolor[rgb]{0.00,0.23,0.31}{\textbf{#1}}}
\newcommand{\NormalTok}[1]{\textcolor[rgb]{0.00,0.23,0.31}{#1}}
\newcommand{\OperatorTok}[1]{\textcolor[rgb]{0.37,0.37,0.37}{#1}}
\newcommand{\OtherTok}[1]{\textcolor[rgb]{0.00,0.23,0.31}{#1}}
\newcommand{\PreprocessorTok}[1]{\textcolor[rgb]{0.68,0.00,0.00}{#1}}
\newcommand{\RegionMarkerTok}[1]{\textcolor[rgb]{0.00,0.23,0.31}{#1}}
\newcommand{\SpecialCharTok}[1]{\textcolor[rgb]{0.37,0.37,0.37}{#1}}
\newcommand{\SpecialStringTok}[1]{\textcolor[rgb]{0.13,0.47,0.30}{#1}}
\newcommand{\StringTok}[1]{\textcolor[rgb]{0.13,0.47,0.30}{#1}}
\newcommand{\VariableTok}[1]{\textcolor[rgb]{0.07,0.07,0.07}{#1}}
\newcommand{\VerbatimStringTok}[1]{\textcolor[rgb]{0.13,0.47,0.30}{#1}}
\newcommand{\WarningTok}[1]{\textcolor[rgb]{0.37,0.37,0.37}{\textit{#1}}}

% ---------- 列表 ----------
\providecommand{\tightlist}{%
  \setlength{\itemsep}{0pt}%
  \setlength{\parskip}{0pt}%
}

% ---------- 表格 ----------
\usepackage{longtable,booktabs,array}
\usepackage{calc}
\usepackage{etoolbox}
\newcounter{none}

\makeatletter
\patchcmd\longtable{\par}{\if@noskipsec\mbox{}\fi\par}{}{}
\makeatother

\IfFileExists{footnotehyper.sty}{\usepackage{footnotehyper}}{\usepackage{footnote}}
\makesavenoteenv{longtable}

% ---------- 图片与 SVG ----------
\usepackage{svg}

\makeatletter
\newsavebox\pandoc@box
\newcommand*\pandocbounded[1]{%
  \sbox\pandoc@box{#1}%
  \Gscale@div\@tempa{\textheight}{\dimexpr\ht\pandoc@box+\dp\pandoc@box\relax}%
  \Gscale@div\@tempb{\linewidth}{\wd\pandoc@box}%
  \ifdim\@tempb\p@<\@tempa\p@\let\@tempa\@tempb\fi
  \ifdim\@tempa\p@<\p@
    \scalebox{\@tempa}{\usebox\pandoc@box}%
  \else
    \usebox\pandoc@box
  \fi
}
\makeatother

% ---------- 链接 ----------
\IfFileExists{xurl.sty}{\usepackage{xurl}}{}
\urlstyle{same}

% ---------- Pandoc 小工具 ----------
\providecommand{\passthrough}[1]{#1}


% ---------- 文档信息 ----------
\newcommand{\BookTitle}{PT-Notes}
\newcommand{\BookAuthor}{Guotao He}
\newcommand{\CoverImage}{titleimage/Ellie.png}

% 日期：yyyy-mm-dd
\newcommand{\twodigits}[1]{\ifnum#1<10 0\fi\number#1}
\renewcommand{\today}{\number\year-\twodigits\month-\twodigits\day}

% ---------- 常用数学宏 ----------
\renewcommand{\vec}[1]{\boldsymbol{#1}}
\newcommand{\ms}[1]{\mathscr{#1}}
\DeclareMathOperator{\sech}{sech}
\DeclareMathOperator{\csch}{csch}
\DeclareMathOperator{\arsinh}{arsinh}
\DeclareMathOperator{\arcosh}{arcosh}
\DeclareMathOperator{\artanh}{artanh}
\let\arSinh\arsinh
\let\arCosh\arcosh
\let\arTanh\artanh

% ---------- 版式 ----------
\linespread{1.3}
\setlength{\lineskip}{2.5pt}
\setlength{\lineskiplimit}{2.5pt}
\setcounter{secnumdepth}{3}
\setcounter{tocdepth}{3}

\renewcommand{\textfraction}{0.15}
\renewcommand{\topfraction}{0.85}
\renewcommand{\bottomfraction}{0.65}
\renewcommand{\floatpagefraction}{0.60}
\floatplacement{figure}{H}
\setkeys{Gin}{width=0.6\textwidth}

% ---------- 标题 ----------
\ctexset{
  chapter = {
    format      = \raggedright,
    number      = \arabic{chapter},
    name        = {Chapter\ ,\ },
    nameformat  = \Large\bfseries,
    aftername   = \par,
    titleformat = \huge\bfseries,
    aftertitle  = {\par\rule[0.5pt]{\textwidth}{1pt}},
    beforeskip  = -10pt,
    afterskip   = 20pt
  },
  part = {
    format      = \centering,
    number      = \Roman{part},
    name        = {Part\ ,\ },
    nameformat  = \Huge\bfseries,
    aftername   = \par,
    titleformat = \Huge\bfseries
  }
}

% ---------- 定理环境 ----------
\newtheoremstyle{boxthmstyle}
  {0pt}{0pt}{\itshape}{}{\bfseries}{.}{0.25em}{}

\mdfdefinestyle{leftbarbox}{
  skipabove=7pt,
  skipbelow=7pt,
  rightline=false,
  topline=false,
  bottomline=false,
  linewidth=4pt,
  innerleftmargin=5pt,
  innerrightmargin=5pt,
  innertopmargin=5pt,
  innerbottommargin=5pt,
  leftmargin=0pt,
  rightmargin=0pt
}

\newmdenv[style=leftbarbox,linecolor=LightSkyBlue,backgroundcolor=LightSkyBlue!15]{thmBox}
\newmdenv[style=leftbarbox,linecolor=Green,backgroundcolor=Green!5]{defiBox}
\newmdenv[style=leftbarbox,linecolor=gray,backgroundcolor=black!5]{lemBox}
\newmdenv[style=leftbarbox,linecolor=TealBlue,backgroundcolor=TealBlue!15]{corollaryBox}
\newmdenv[style=leftbarbox,linecolor=red,backgroundcolor=red!5]{attentionBox}

\theoremstyle{boxthmstyle}
\newtheorem{thm}{Theorem}[section]
\newtheorem{ppt}[thm]{Proposition}
\newtheorem{lem}[thm]{Lemma}
\newtheorem{mycorollary}[thm]{Corollary}
\newtheorem{defi}{Definition}[section]

\newenvironment{theorem}{\begin{thmBox}\begin{thm}}{\end{thm}\end{thmBox}}
\newenvironment{proposition}{\begin{thmBox}\begin{ppt}}{\end{ppt}\end{thmBox}}
\newenvironment{lemma}{\begin{lemBox}\begin{lem}}{\end{lem}\end{lemBox}}
\newenvironment{corollary}{\begin{corollaryBox}\begin{mycorollary}}{\end{mycorollary}\end{corollaryBox}}
\newenvironment{definition}{\begin{defiBox}\begin{defi}}{\end{defi}\end{defiBox}}

\theoremstyle{definition}
\newtheorem{example}{Example}[section]
\newtheorem{question}{Question}[section]
\newtheorem{remark}{Remark}[section]
\newenvironment{solution}{\begin{proof}[\bfseries Solution]}{\end{proof}}
\renewcommand{\proofname}{Proof}
\newenvironment{attention}{\begin{attentionBox}\textbf{Attention.}\ }{\end{attentionBox}}

% ---------- PDF 链接 ----------
\hypersetup{
  unicode=true,
  colorlinks=true,
  linkcolor=MidnightBlue,
  urlcolor=TealBlue,
  citecolor=OliveGreen,
  pdftitle={\BookTitle},
  pdfauthor={\BookAuthor}
}

\begin{document}

\begingroup
\newgeometry{margin=0pt}
\thispagestyle{empty}

\centering
\IfFileExists{\CoverImage}{%
  \includegraphics[
    width=\paperwidth,
    height=0.55\paperheight,
    keepaspectratio
  ]{\CoverImage}%
}{%
  \vspace*{0.18\paperheight}%
}

\vfill
{\parbox{0.618\textwidth}{%
  \raggedleft
  {\bfseries\Huge \BookTitle\par}
  \vspace{0.6em}
  \rule{\linewidth}{4pt}
}}

\vfill
{\parbox{0.618\textwidth}{%
  \raggedleft\Large\kaishu
  \begin{tabular}{r|}
    作者：\BookAuthor \\
    编译时间：\today
  \end{tabular}
}}
\vfill

\restoregeometry
\endgroup

\clearpage

\frontmatter
\thispagestyle{empty}
\vspace*{\fill}

\noindent
本文档对应 PT-Notes 项目，最新内容见
\url{https://github.com/GHe0000/PT-Notes}。建议通过项目主页或仓库分享，以便获得最新版本。\par\medskip

\noindent
本文档整理了 PT 比赛中使用和记录的实验方法、工具与经验。内容以实践参考为目的，不保证所有表述在所有条件下均严格适用；如发现错误或不清晰之处，欢迎通过仓库 Issue 或 Pull Request 反馈。\par\medskip

\noindent
正文中引用的书籍、论文、网页、图片和其他第三方材料，其版权归原作者所有。除另有说明外，本文档正文按照 CC BY-NC 4.0 许可协议共享，完整许可证见
\url{https://creativecommons.org/licenses/by-nc/4.0}。\par\medskip

\noindent
Copyright \copyright\ 2026--\the\year\ \BookAuthor\\
\textit{Not Published}

\clearpage
\tableofcontents
\clearpage

% index.qmd（前言）仍处于 frontmatter；
% index.qmd 末尾的 Raw LaTeX 会执行 \mainmatter，
% 因此“常用软件”从 Chapter 1 开始编号。
\IfFileExists{generated/book-latex/PT-Notes.tex}{%
  \subimport{generated/book-latex/}{PT-Notes.tex}%
}{%
  \IfFileExists{generated/PT-Notes.tex}{%
    \subimport{generated/}{PT-Notes.tex}%
  }{%
    \PackageWarningNoLine{PT-Notes}{Quarto-generated PT-Notes.tex not found; compiling Chapter/test.tex instead}%
    \mainmatter
    % !TeX root = ../main.tex
\chapter{环境测试}

\section{正文与数学}

这是一段中文正文，用于检查 XeLaTeX 中文排版、行距、目录与章节标题。
行内公式示例：$\vec{x}\in\mathbb{R}^n$，以及花体符号 $\ms{F}$。

\begin{align}
  \mathrm{e}^{\mathrm{i}x} &= \cos x + \mathrm{i}\sin x, \\
  \int_{-\infty}^{+\infty} \mathrm{e}^{-x^2}\,\mathrm{d}x &= \sqrt{\pi}.
\end{align}

旧式 Dirac 记号兼容测试：
\[
  \bra{\psi},\qquad
  \ket{\phi},\qquad
  \braket{\psi|\phi},\qquad
  \braket{\psi|A|\phi}.
\]

反双曲函数宏：$\arsinh x$，$\arcosh x$，$\artanh x$；单位示例：\SI{9.81}{\metre\per\second\squared}。

\section{定理类环境}

\begin{definition}[测试定义]\label{def:test}
设 $V$ 为线性空间。这是 Definition 环境。
\end{definition}

\begin{theorem}[测试定理]\label{thm:test}
若 $a,b\in\mathbb{R}$，则 $(a+b)^2=a^2+2ab+b^2$。
\end{theorem}

\begin{proof}
直接展开右端即可。
\end{proof}

\begin{proposition}[测试命题]
这是 Proposition 环境，并与 Theorem 共用计数器。
\end{proposition}

\begin{lemma}[测试引理]
这是 Lemma 环境。
\end{lemma}

\begin{corollary}[测试推论]
这是 Corollary 环境。
\end{corollary}

\begin{example}[测试例子]
由定理~\ref{thm:test}，取 $a=b=1$ 即得 $4=4$。
\end{example}

\begin{question}[测试问题]
Definition~\ref{def:test} 的编号和 Theorem 是否独立？
\end{question}

\begin{solution}
是。Definition 使用独立计数器。
\end{solution}

\begin{remark}[测试注记]
Remark 使用普通的 definition theorem style，不带彩色边框。
\end{remark}

\begin{attention}
这是 Attention 环境，用于提示需要特别注意的内容。
\end{attention}

\section{Quarto 兼容环境}

\begin{itemize}
  \tightlist
  \item Quarto/Pandoc 紧凑列表测试；
  \item 第二项。
\end{itemize}

\begin{Shaded}
\begin{Highlighting}[]
\NormalTok{print}\OperatorTok{(}\StringTok{"hello"}\OperatorTok{)}
\end{Highlighting}
\end{Shaded}

\section{列表、表格与图片}

\begin{table}[H]
  \centering
  \caption{表格环境测试}
  \begin{tabular}{c c}
    \toprule
    参数 & 数值 \\
    \midrule
    $a$ & 1 \\
    $b$ & 2 \\
    \bottomrule
  \end{tabular}
\end{table}

\begin{figure}[H]
  \centering
  \rule{0.6\textwidth}{0.28\textwidth}
  \caption{图片环境占位测试}
\end{figure}

\section{链接}

内部引用：见定理~\ref{thm:test}。外部链接测试：\url{https://www.latex-project.org/}。

  }%
}

\end{document}
